\section{Reception: The Master Nobody Canonized}

\subsection{Ancient and late-antique judgment}

The shape of Tertullian's reception was fixed early: indispensable and
suspect, in that order. Cyprian, bishop of Carthage within a generation of
him, is reported never to have passed a day without reading him, asking his
secretary, ``Give me the master'' (Jerome, \work{On Illustrious Men} 53)---and
Cyprian's surviving treatises rework Tertullian's themes without naming him. Lactantius, the first to survey Latin Christian letters, granted that
Tertullian ``was skilled in literature of every kind'' but complained that ``in
eloquence he had little readiness, and was not sufficiently polished, and very
obscure'' (\work{Divine Institutes} 5.1)---a stylist's verdict on a style that
the following centuries nevertheless kept quoting. Eusebius cited the
\work{Apology} with respect; Jerome made him ``chief of the Latin writers''
after none.

Augustine's engagement is the most instructive, because it runs in both
directions at once. Cataloguing heresies in 428, he made Tertullian a heretic
and the Tertullianists a sect (\work{On Heresies} 86); answering a
speculative layman on the soul's origin, he congratulated the man for having
``kept himself uninfluenced by the ravings of Tertullian,'' who ``insisted,
that as the soul is corporeal, so likewise is God'' (\work{On the Soul and
its Origin} 2.9). Yet the same Augustine demonstrably worked with
Tertullian's \work{A Treatise on the Soul} and the moral treatises: the
African master was refuted by name and used without it. Vincent of L\'erins
then gave the ambivalence its classical form: ``who more learned than he, who
more versed in knowledge whether divine or human?''---and yet, adopting
Hilary's judgment that ``by his subsequent error he detracted from the
authority of his approved writings,'' Vincent concluded that Tertullian, like
Origen among the Greeks, ``was a great trial in the Church''
(\work{Commonitorium} 18). By the sixth century the so-called Gelasian
Decree, a Roman index of receivable and non-receivable books, simply listed
the works of Tertullian among the \textit{apocrypha}, between the history of
Eusebius and the works of Lactantius---not condemned for a doctrine, but not
to be received as authoritative either (\work{Decretum Gelasianum}, list of
apocrypha).

\subsection{Survival by a thread: the manuscripts}

What a church will not canonize it may still copy, and the copying is its own
verdict: grudging, thin, and just sufficient. The \work{Apology} travelled
well---at least thirty-nine manuscripts, many containing it alone. The rest
of the corpus survived in a handful of late collections, of which the
specialist census (this study follows the Tertullian Project's) counts five:
the great Cluny corpus of twenty-odd works, copied from the eleventh century
onward; the Agobardine corpus, whose one surviving witness, Codex
Agobardinus (Paris, Biblioth\`eque nationale, lat.~1622), a ninth-century
book once owned by Bishop Agobard of Lyons, has lost its back half and
preserves thirteen works of the twenty-one its own table of contents
promises; the Trecensis corpus, surviving in a single twelfth-century
manuscript; the Ottobonian corpus, four works in one Vatican manuscript
identified only in the twentieth century; and the Corbie corpus, which
gathered precisely the works reckoned heretical, and every manuscript of
which had disappeared by about 1600. The loss fell exactly where the corpus
was most inflammatory, and it means that several treatises---\work{On
Fasting} and \work{On Modesty} among them---survive today only because
printers reached them first: Beatus Rhenanus produced the first printed
edition at Basel in 1521 from two manuscripts, one now lost, and the editions
of Mesnart (1545) and Gelenius (1550) printed works whose manuscript sources
subsequently vanished. A substantial part of what this study quotes therefore rests on
sixteenth-century printings of books no one has seen since. The margin by
which the first Latin theologian escaped oblivion was, at several points, one
copy wide.

\subsection{Not a saint, and yet a Father}

Tertullian was never canonized and has no feast; no cult, calendar entry, or
liturgical commemoration of him, ancient or modern, was located in the
research for this study (a bounded negative search; see the appendix). He is likewise no
Doctor of the Church---a title requiring sanctity as well as eminent doctrine.
Yet he is universally treated as a Church Father in the broad sense: the
patristic manuals, the critical editions of the Fathers, and the Church's own
documents handle him as a primary witness to ancient faith and practice. The
\textit{Catechism of the Catholic Church} quotes his \work{On Prayer} for the
Lord's Prayer as ``the summary of the whole Gospel'' (CCC 2761); Benedict XVI
gave him an entire general audience in his patristic catecheses (30 May 2007),
presenting the man who began Christian Latin literature, crediting him
expressly with introducing the Latin terms ``one substance'' and ``three
Persons'' for the Trinitarian mystery, and quoting the seed-of-blood epigram
back into the Church's ordinary teaching---and who yet, in the same
catechesis, stands as the great theologian who lacked the humility to remain
with the Church and to bear his own and others' weaknesses; for Benedict, a
standing lesson that theological brilliance without that humility consumes
itself. Selective citation, it
should be stressed, is reception of particular formulations, not
rehabilitation of the whole corpus or a judgment on the man's end.

\subsection{The contested pages}

An honest account keeps the pages later ages winced at. The treatise \work{On
the Apparel of Women} addresses its readers, on the strength of Eve's
primacy in the Fall, in the sentence for which Tertullian is now most often
indicted: ``You are the devil's gateway: you are the unsealer of that
(forbidden) tree: you are the first deserter of the divine law'' (\work{On the
Apparel of Women} 1.1). It stands in a genre of moral invective against luxury,
in a corpus that also praises the shared discipleship of husband and wife
(\work{To His Wife} 2.8), honors Perpetua as ``the most heroic martyr''
(\work{On the Soul} 55), and reports women exercising recognized prophetic
gifts in the liturgy (\work{On the Soul} 9); but the sentence is his, the
theology of blame in it is his, and no genre plea dissolves it. Similarly, the
closing chapter of \work{The Shows} savors the damnation of the persecutors,
philosophers, and actors as the best spectacle of all---``What there excites my
admiration? What my derision? Which sight gives me joy?'' (\work{The Shows}
30)---a page that has served every critic of Christian vindictiveness from the
Enlightenment onward. The rigorist imagination that made Tertullian a matchless
witness against comfortable Christianity is the same one that could not wait
for mercy. Both facts belong to his legacy.

A third contested page exists only in reception. The slogan ``credo quia
absurdum,'' endlessly attributed to him as the motto of fideist unreason,
occurs nowhere in his works; what he wrote (\work{On the Flesh of Christ} 5,
quoted with its Latin in the theology section) is a controlled paradox about
why so shameful a story would never have been invented. The misquotation has
functioned for centuries as a portable indictment of Christian
irrationality, and its career tells more about his reputation than about his
text: Tertullian became the man to whom Western culture ascribes the sayings
it needs an extremist to have said. The remedy, here as throughout, is the
locus itself.

\subsection{The modern rebalancing}

Modern scholarship, meanwhile, quietly returned him to the Church's history,
and the return can be dated by its instruments. The Vienna corpus (CSEL)
issued critical texts of Tertullian in successive volumes from 1890 to 1957;
the twentieth century added the studies that dismantled the legendary
biography (Barnes 1971) and re-examined the ``schism'' (Powell 1975; Rankin
1995). Eric Osborn's \work{Tertullian, First Theologian of the West}
(Cambridge, 1997) argued---in its publisher's summary, for the book itself
was not read for this study---that a thinker ``for a long time misquoted and
misused'' now calls for sustained analysis, and offered a full reappraisal of
his theology's influence on the Western tradition; Geoffrey Dunn's
\work{Tertullian} (Routledge, 2004) put fresh translations and an accessible
introduction into student hands. Where the nineteenth century wrote
confidently of ``Tertullian the Montanist,'' the current literature writes of
a Catholic controversialist of rigorist temper whose New Prophecy allegiance
need not have carried him out of communion, whose ``sect'' may be an artifact
of fifth-century cataloguing, and whose lay voice---if lay he was---is among
the most important we possess for ante-Nicene Christianity. The paradox
Vincent named remains: the Latin Church's first master theologian is the one
it can quote but not celebrate.
